\section{Czujnik klimatyczny}
Centralnym elementem warstwy odpowiedzialnej za pomiary środowiskowe 
jest struktura \texttt{ClimateSensor}, przedstawiona na 
listingu~\ref{lst:climate_sensor_struct}, której zadaniem jest agregacja 
trzech niezależnych modułów pomiarowych — czujnika temperatury 
i~wilgotności AHT20, czujnika ciśnienia i~temperatury BMP280 oraz 
czujnika natężenia oświetlenia VEML7700 — pod wspólnym, 
jednolitym interfejsem. 
Struktura ta hermetyzuje nie tylko same sterowniki poszczególnych 
czujników, lecz również współdzieloną magistralę komunikacyjną \isqrc{}, 
mechanizm opóźnień (\texttt{DelayNs}) oraz piny przetwornika ADC 
wykorzystywane do pomiaru napięcia na kondensatorach układu zasilania. 
Dzięki temu logika odpowiedzialna za obsługę pojedynczego węzła pomiarowego 
została skoncentrowana w~jednym miejscu, co znacząco upraszcza jej dalsze 
wykorzystanie w~warstwie wyższego poziomu, odpowiedzialnej za transmisję danych.

\begin{lstlisting}[
    language=Rust, style=colouredRust, 
    caption={Struktura \texttt{ClimateSensor}},
    label={lst:climate_sensor_struct}
]
pub struct ClimateSensor<
    VemlConfig, 
    Bmp280Config, 
    D: DelayNs,
    SumCapPin: AdcChannel<
        arduino_hal::hal::Atmega, 
        arduino_hal::pac::ADC>,
    SecondCapPin: AdcChannel<
        arduino_hal::hal::Atmega, 
        arduino_hal::pac::ADC>,
>{
    // SENSORS
    aht20: Aht20,
    bmp280: Bmp280<Bmp280Config>,
    veml7700: Veml7700<VemlConfig>,
    // PINS
    sum_capacitor: SumCapPin,
    capacitor_2_pin: SecondCapPin,
    // COMMUNICATION
    i2c: i2c::I2c,
    // TIMING
    delay: D,
    // MISC
    sensor_id: u8,
}
\end{lstlisting}

Parametryzacja generyczna struktury (\texttt{VemlConfig}, \texttt{Bmp280Config}, 
\texttt{D}, \texttt{SumCapPin}, \texttt{SecondCapPin}) pozwala na statyczne, 
rozstrzygane na etapie kompilacji dostosowanie konfiguracji poszczególnych 
modułów oraz przypisanych pinów, bez wprowadzania narzutu związanego 
z~dynamicznym dyspozytorem \translation{dynamic dispatch}.

Za centralizację odczytu odpowiadają metody \texttt{read\_modules} oraz 
\texttt{read\_charge\_info}, które sekwencyjnie odpytują poszczególne czujniki, 
a~uzyskane wartości kodują przy użyciu typu \texttt{TypedLabelReadout}, 
zapisując je bezpośrednio do wspólnego bufora bajtowego. 
Metoda \texttt{read\_bytes}, stanowiąca publiczny punkt wejścia dla warstw 
nadrzędnych, łączy identyfikator węzła pomiarowego, dane z~modułów klimatycznych 
oraz informacje o~stanie naładowania kondensatorów w~jedną, spójną ramkę danych 
o~ustalonym rozmiarze, gotową do dalszej transmisji.

Obsługa błędów została zorganizowana w~sposób hierarchiczny, poprzez zestaw 
dedykowanych typów wyliczeniowych (\texttt{ModuleInitError}, 
\texttt{ClimateSensorInitError}, \texttt{ModuleReadError}, 
\texttt{ClimateSensorReadError}) wraz z~implementacjami cechy \texttt{From}, 
umożliwiającymi automatyczną propagację błędów niższego poziomu za 
pomocą operatora \texttt{?}.